% problem_id: S001-lemma-backwards-homotopy
% source_id: S001 (Stacks Project)
% source_file: simplicial.tex
% statement_locator: simplicial.tex lines 5289--5302
% proof_locator: simplicial.tex lines 5304--5500
% env: lemma
% id_note: label 在本来源内唯一
% pairing: tight (verified)
% status: complete (statement + author proof verbatim + reason + locators)
%
\begin{lemma}
\label{lemma-backwards-homotopy}
Let $\mathcal{A}$ be an abelian category.
Let $U$, $V$ be simplicial objects of $\mathcal{A}$.
Let $a, b : U \to V$ be a pair of morphisms.
Assume the corresponding maps of chain complexes
$N(a), N(b) : N(U) \to N(V)$ are homotopic by
a homotopy $\{N_n : N(U)_n \to N(V)_{n + 1}\}$.
Then there exists a homotopy from $a$ to $b$ as in
Definition \ref{definition-homotopy}. Moreover, one can choose the
homotopy $h : U \times \Delta[1] \to V$ such that
$N_n = N(h)_n$ where $N(h)$ is the homotopy coming
from $h$ as in Section \ref{section-homotopy-abelian}.
\end{lemma}

\begin{proof}
Let $(\diamond N(U), \diamond a, \diamond b, \diamond h)$
be as in Lemma \ref{lemma-represent-homotopy} and its proof. By that lemma
there exists a morphism $\diamond N(U) \to N(V)$ representing
the triple $(N(a), N(b), \{N_n\})$. We will show
there exists a morphism
$\psi : N(U \times \Delta[1]) \to \diamond{N(U)}$
such that $\diamond a = \psi \circ N(e_0)$, and
$\diamond b = \psi \circ N(e_1)$. Moreover, we will
show that the homotopy between $N(e_0), N(e_1) :
N(U) \to N(U \times \Delta[1])$ coming from
(\ref{equation-homotopy-to-homotopy}) and
Lemma \ref{lemma-homotopy-s-N} with
$h = \text{id}_{U \times \Delta[1]}$ is mapped via
$\psi$ to the canonical homotopy $\diamond h$ between the two
maps $\diamond a, \diamond b : N(U) \to \diamond{N(U)}$. Certainly this
will imply the lemma.

\medskip\noindent
Note that $N : \text{Simp}(\mathcal{A}) \to \text{Ch}_{\geq 0}(\mathcal{A})$
as a functor is a direct summand of the functor
$s : \text{Simp}(\mathcal{A}) \to \text{Ch}_{\geq 0}(\mathcal{A})$.
Also, the functor $\diamond$ is compatible with direct sums.
Thus it suffices instead to construct a morphism
$\Psi : s(U \times \Delta[1]) \to \diamond{s(U)}$ with the
corresponding properties. This is what we do below.

\medskip\noindent
By Definition \ref{definition-homotopy}
the morphisms $e_0 : U \to U \times \Delta[1]$
and $e_1 : U \to U \times \Delta[1]$ are homotopic
with homotopy $\text{id}_{U \times \Delta[1]}$.
By Lemma \ref{lemma-homotopy-s-N}
we get an explicit homotopy
$\{h_n : s(U)_n \to s(U \times \Delta[1])_{n + 1}\}$
between the morphisms
of chain complexes $s(e_0) : s(U) \to s(U \times \Delta[1])$
and $s(e_1) : s(U) \to s(U \times \Delta[1])$. By
Lemma \ref{lemma-map-into-diamond} above we get a corresponding morphism
$$
\Phi : \diamond{s(U)} \to s(U \times \Delta[1])
$$
According to the construction, $\Phi_n$ restricted to the summand
$s(U)[-1]_n = s(U)_{n - 1}$ of $\diamond{s(U)}_n$
is equal to $h_{n - 1}$. And
$$
h_{n - 1} = \sum\nolimits_{i = 0}^{n - 1}
(-1)^{i + 1} s^n_i \cdot \alpha^n_{i + 1} :
U_{n - 1} \to \bigoplus\nolimits_j U_n \cdot \alpha^n_j.
$$
with obvious notation.

\medskip\noindent
On the other hand, the morphisms $e_i : U \to U \times \Delta[1]$ induce
a morphism $(e_0, e_1) : U \oplus U \to U \times \Delta[1]$.
Denote $W$ the cokernel. Note that, if we write
$(U \times \Delta[1])_n = \bigoplus_{\alpha : [n] \to [1]} U_n \cdot \alpha$,
then we may identify
$W_n = \bigoplus_{i = 1}^n U_n \cdot \alpha^n_i$ with
$\alpha^n_i$ as in Section \ref{section-homotopy}.
We have a commutative diagram
$$
\xymatrix{
0 \ar[r] &
U \oplus U \ar[rd]_{(1, 1)} \ar[r] &
U \times \Delta[1] \ar[d]^\pi \ar[r] &
W \ar[r] &
0 \\
& & U & &
}
$$
This implies we have a similar commutative diagram after
applying the functor $s$. Next, we choose the splittings
$\sigma_n : s(W)_n \to s(U \times \Delta[1])_n$ by mapping
the summand $U_n \cdot \alpha^n_i \subset W_n$ via $(-1, 1)$
to the summands $U_n \cdot \alpha^n_0 \oplus U_n \cdot \alpha^n_i
\subset (U \times \Delta[1])_n$. Note that $s(\pi)_n \circ \sigma_n = 0$.
It follows that $(1, 1) \circ \delta(\sigma)_n = 0$.
Hence $\delta(\sigma)$ factors as in
Lemma \ref{lemma-map-into-diamond}. By that lemma
we obtain a canonical morphism
$\Psi : s(U \times \Delta[1]) \to \diamond{s(U)}$.

\medskip\noindent
To compute $\Psi$ we first compute the morphism
$\delta(\sigma) : s(W) \to s(U)[-1] \oplus s(U)[-1]$.
According to Homology, Lemma \ref{homology-lemma-ses-termwise-split}
and its proof,
to do this we have compute
$$
d_{s(U \times \delta[1]), n} \circ \sigma_n
-
\sigma_{n - 1} \circ d_{s(W), n}
$$
and write it as a morphism into
$U_{n - 1} \cdot \alpha^{n - 1}_0 \oplus U_{n - 1} \cdot \alpha^{n - 1}_n$.
We only do this in case $\mathcal{A}$ is the category of abelian
groups. We use the short hand notation $x_{\alpha}$ for $x \in U_n$
to denote the element $x$ in the summand $U_n \cdot \alpha$
of $(U \times \Delta[1])_n$. Recall that
$$
d_{s(U \times \delta[1]), n}
=
\sum\nolimits_{i = 0}^n (-1)^i d^n_i
$$
where $d^n_i$ maps the summand $U_n \cdot \alpha$
to the summand $U_{n - 1} \cdot (\alpha \circ \delta^n_i)$
via the morphism $d^n_i$ of the simplicial object $U$.
In terms of the notation above this means
$$
d_{s(U \times \delta[1]), n}(x_\alpha) =
\sum\nolimits_{i = 0}^n (-1)^i (d^n_i(x))_{\alpha \circ \delta^n_i}
$$
Starting with $x_\alpha \in W_n$, in other words $\alpha = \alpha^n_j$
for some $j \in \{1, \ldots, n\}$, we see that
$\sigma_n(x_\alpha) = x_\alpha - x_{\alpha^n_0}$ and
hence
$$
(d_{s(U \times \delta[1]), n} \circ \sigma_n)(x_\alpha)
=
\sum\nolimits_{i = 0}^n (-1)^i (d^n_i(x))_{\alpha \circ \delta^n_i}
-
\sum\nolimits_{i = 0}^n (-1)^i (d^n_i(x))_{\alpha^n_0 \circ \delta^n_i}
$$
To compute $d_{s(W), n}(x_\alpha)$, we have to omit
all terms where $\alpha \circ \delta^n_i = \alpha^{n - 1}_0, \alpha^{n - 1}_n$.
Hence we get
$$
\begin{matrix}
(\sigma_{n - 1} \circ d_{s(W), n})(x_\alpha) = \\
\sum\nolimits_{i = 0, \ldots, n\text{ and }
\alpha \circ \delta^n_i \not = \alpha^{n - 1}_0\text{ or }\alpha^{n - 1}_n}
\Big((-1)^i (d^n_i(x))_{\alpha \circ \delta^n_i}
-
(-1)^i (d^n_i(x))_{\alpha^{n - 1}_0}
\Big)
\end{matrix}
$$
Clearly the difference of the two terms is the sum
$$
\sum\nolimits_{i = 0, \ldots, n\text{ and }
\alpha \circ \delta^n_i = \alpha^{n - 1}_0\text{ or }\alpha^{n - 1}_n}
\Big((-1)^i (d^n_i(x))_{\alpha \circ \delta^n_i}
-
(-1)^i (d^n_i(x))_{\alpha^{n - 1}_0}
\Big)
$$
Of course, if $\alpha \circ \delta^n_i = \alpha^{n - 1}_0$
then the term drops out. Recall that $\alpha = \alpha^n_j$
for some $j \in \{1, \ldots, n\}$. The only way
$\alpha^n_j \circ \delta^n_i = \alpha^{n - 1}_n$
is if $j = n$ and $i = n$. Thus we actually get
$0$ unless $j = n$ and in that case we get
$(-1)^n(d^n_n(x))_{\alpha^{n - 1}_n} - (-1)^n(d^n_n(x))_{\alpha^{n - 1}_0}$.
In other words, we conclude the morphism
$$
\delta(\sigma)_n :
W_n
\to
(s(U)[-1] \oplus s(U)[-1])_n = U_{n - 1} \oplus U_{n - 1}
$$
is zero on all summands except $U_n \cdot \alpha^n_n$
and on that summand it is equal to $((-1)^nd^n_n, -(-1)^nd^n_n)$.
(Namely, the first summand of the two corresponds to the factor
with $\alpha^{n - 1}_n$ because that is the map $[n - 1] \to [1]$
which maps everybody to $0$, and hence corresponds to $e_0$.)

\medskip\noindent
We obtain a canonical diagram
$$
\xymatrix{
0 \ar[r] &
s(U) \oplus s(U) \ar[r] \ar[d] &
\diamond{s(U)} \ar[r] \ar[d]^{\Phi}&
s(U)[-1] \ar[r] \ar[d] &
0 \\
0 \ar[r] &
s(U) \oplus s(U) \ar[r] \ar[d] &
s(U \times \Delta[1]) \ar[r] \ar[d]^\Psi &
s(W) \ar[r] \ar[d] &
0 \\
0 \ar[r] &
s(U) \oplus s(U) \ar[r] &
\diamond{s(U)} \ar[r] &
s(U)[-1] \ar[r] &
0
}
$$
We claim that $\Phi \circ \Psi$ is the identity.
To see this it is enough to prove that the composition
of $\Phi$ and $\delta(\sigma)$ as a map
$s(U)[-1] \to s(W) \to s(U)[-1] \oplus s(U)[-1]$ is the
identity in the first factor and minus identity in the second.
By the computations above it is
$((-1)^nd^n_0, -(-1)^nd^n_0) \circ (-1)^n s^n_n = (1, -1)$
as desired.
\end{proof}
